\documentclass[10pt,a4paper]{amsart}
\usepackage[margin=19mm]{geometry}
\usepackage{amsmath,amssymb,mathtools}
\usepackage[scr=rsfs,cal=euler]{mathalfa}
\usepackage{xfrac}
\usepackage[unicode,hidelinks,bookmarks=false]{hyperref}
\newcommand{\ie}{i.e.\ }
\newcommand{\cf}{cf.\ }
\newcommand{\resp}{respectively\ }
\newcommand{\Subsection}{\subsection}
\providecommand{\nobreakdash}{\nobreak\hbox{-}}
\setlength{\emergencystretch}{2em}
\multlinegap0pt
\usepackage{bm}
\usepackage{bbm}
\usepackage{dsfont}
\usepackage{xspace}
\usepackage{cancel}
\usepackage{mathtools}
\usepackage{amsthm}
\makeatletter
\@ifundefined{Thm}{\newtheorem{Thm}{Thm}}{}
\@ifundefined{Lem}{\newtheorem{Lem}[Thm]{Lemma}}{}
\makeatother
\newcommand{\R}{\mathbb{R}}
\title[Blowing-up solutions to a critical 4D Neumann system in a co]{Blowing-up solutions to a critical 4D Neumann system in a competitive regime}
\author{Qing Guo, Angela Pistoia, Shixin Wen}
\date{Source: arXiv:2603.29329v1 (CC BY 4.0)}
\begin{document}
\maketitle

\noindent\emph{Source and licence.} Blowing-up solutions to a critical 4D Neumann system in a competitive regime. Version arXiv:2603.29329v1 (CC BY 4.0), distributed under the Creative Commons Attribution 4.0 International licence, as recorded in the arXiv licence metadata for this exact version and shown on the abstract page https://arxiv.org/abs/2603.29329v1. Full rights notice: licenses/arxiv-2603.29329-CC-BY-4.0.txt.\par\smallskip

\noindent\textit{Editorial scope.} Counted unit: the Lem whose source label is \texttt{invert} in the article (source0.tex lines 512--517, source label \texttt{invert}), delivered together with the article's own complete original proof of it (source0.tex lines 518--694, one \texttt{proof} environment of 177 lines). No mathematics is written, repaired or re-derived: statement and proof are the article's own text line for line. Required context retained uncounted: vec-1 (lines 332--335); properties-1 (lines 479--484). Every definition, display and label of those lines is retained unchanged. Omitted from this package and not redistributed: 1--331, 336--478, 485--511, 695--1019. Nothing inside a retained excerpt is omitted or altered: every retained body is delivered exactly as the article's lines are written, so no omission receipt of any kind accompanies this package. Citations: the retained excerpts carry no citation command at all, so no bibliography entry of the article is transcribed and no reference is redistributed. Reference adaptation: none was needed and none was made; every reference inside the delivered unit resolves inside it, so no unresolved reference or citation is produced. Printed slips kept literal: no printed word, symbol, hypothesis or number of the counted statement or of its proof was altered. Layout and class: the article's own class is \texttt{amsart} with packages including tikz, tikz-3dplot, inputenc, fontenc, amsmath, amssymb, amsfonts, amsthm, geometry, graphicx, mathrsfs, latexsym, color, hyperref; the wrapper is the lane's pinned \texttt{amsart} editorial wrapper at 10pt on A4, which restates the theorem-like environments the retained text needs and adds the article's own macro definitions that the retained text needs (\texttt{\textbackslash R}), with extra wrapper packages (bm, bbm, dsfont, xspace, cancel, mathtools). The delivered numbering is the unit's own. Assets: no retained span contains a figure environment, a graphics-inclusion command, a PDF-page inclusion, a table or a TikZ picture; no image file is copied into this package, the package declares no asset dependency and no image file is redistributed with it. Licence: version arXiv:2603.29329v1 of this article is distributed under the Creative Commons Attribution 4.0 International licence, as recorded in the arXiv licence metadata for this exact version and shown on the abstract page https://arxiv.org/abs/2603.29329v1. A full rights notice is delivered at licenses/arxiv-2603.29329-CC-BY-4.0.txt. Characters: every retained excerpt is pure ASCII, so the delivered engine, XeLaTeX with its default Latin Modern text font, reports no missing character.
 \begin{align}
    \label{vec-1} \Pi^{\perp}_{\bm{d},\bm{\xi}}\left(\bm{\mathcal{L}}(\bm{\psi})-\bm{\mathcal{E}}-\bm{\mathcal{N}}(\bm{\psi})\right)&=0\\
    \label{vec-2}\Pi_{\bm{d},\bm{\xi}}\left(\bm{\mathcal{L}}(\bm{\psi})-\bm{\mathcal{E}}-\bm{\mathcal{N}}(\bm{\psi})\right)&=0.
\end{align}

    \begin{equation}\label{properties-1}
    \begin{cases}
        W(r)\lesssim |\ln{r}|, &r\to 0.\\
        W(r)\lesssim \frac{1}{r^{2}},&r\to +\infty.
    \end{cases}
    \end{equation}

 {\begin{Lem}\label{invert}
 Given any $\eta\in(0,1)$ and $b>0$, there exists $c>0$ and $\lambda_0>0$ such that for any $\lambda \in (\lambda_0,+\infty)$, for any $\left(\bm{d}, \bm{\xi}\right)\in \mathcal{O}_\eta$ and for any $\beta\in (0,b)$ 
    \begin{align*}
        \left\|\left(\Pi^{\perp}_{\bm{d},\bm{\xi}}\circ\bm{\mathcal{L}}_{\bm{d},\bm{\xi}}\right)(\bm{\psi})\right\|_H\geq c\left\|\bm{\psi}\right\|_H, \quad \forall \bm{\psi}\in \bm{K^{\perp}}_{\bm{d},\bm{\xi}}.   \end{align*} 

\end{Lem}}

\begin{proof} Suppose by contradiction that, there exist sequences $\bm{\psi_{{n}}} \in \bm{K^{\perp}}_{\bm{d},\bm{\xi}}$ with $\left\|{\bm{\psi_{{n}}}}\right\|_{H}=1$, and $\bm{h_n}\in\bm{K^{\perp}}_{\bm{d},\bm{\xi}}$ such that   
    \begin{align*}
\left\|\bm\Pi^{\perp}\circ\bm{\mathcal{L}}(\bm{\psi_n})\right\|_{H}&=\left\|\bm{{h}_{n}}\right\|_{H} \rightarrow 0,\quad\text{as $n\rightarrow +\infty$}.
    \end{align*}
From \eqref{vec-1}, we derive that for $i,j=1,2$ with $i\neq j$,
\begin{align}\label{new-psi}
   \psi_{in}-i_{\lambda_n}^{*}\left(3V_{in}^2\psi_{in}-\beta\left(V_{jn}^2\psi_{in}+2V_{in}V_{jn}\psi_{jn}\right)\right)=h_{in}+w_{in},
\end{align}
where $w_{{in}} \in K_i$.


\medskip
\textbf{Step 1.}~ We prove that $\bm{w_{{n}}} \rightarrow 0$ strongly in $H$ as $n \rightarrow+\infty$. For this aim, we need to prove that for $i=1,2$, $w_{in}\rightarrow 0$  strongly in $H^1(\Omega)$ as $n \rightarrow+\infty$. 

Indeed, since $w_{in}\in K_i$, then there exist $\{c_{in}^{(k)}\}_{k=0}^{3}$, such that 
\begin{align}
    \label{w-1}\|w_{in}\|^2=\sum_{k=0}^3\left(c_{in}^{(k)}\right)^2\sigma_{kk}+\sum_{k,j=0,\atop k\neq j}^3\left(c_{in}^{(k)}\,c_{in}^{(j)}\right)\sigma_{kj},
\end{align}
where
\begin{align}\label{useful}
\sigma_{00}:=&\displaystyle\int_{\Omega}\left|\nabla\left(\delta_{in}\,\frac{\partial U_{in}}{\partial \delta_i}\right)\right|^2+\lambda_n\displaystyle\int_{\Omega}\left|\delta_{in}\,\frac{\partial U_{in}}{\partial \delta_i}\right|^2=\widetilde{A_0}+O(\delta_{in}),\nonumber\\\sigma_{kk}:=&\displaystyle\int_{\Omega}\left|\nabla\left(\delta_{in}\,\frac{\partial U_{in}}{\partial t_{k,i}}\right)\right|^2+\lambda_n\displaystyle\int_{\Omega}\left|\delta_{in}\,\frac{\partial U_{in}}{\partial t_{k,i}}\right|^2=\widetilde{A_1}+O(\delta_{in}),\nonumber\\
\sigma_{0j}:=&\displaystyle\int_{\Omega}\nabla\left(\delta_{in}\,\frac{\partial U_{in}}{\partial\delta_{i}}\right)\nabla\left(\delta_{in}\,\frac{\partial U_{in}}{\partial t_{j,i}}\right)+\lambda_n\displaystyle\int_{\Omega}\left(\delta_{in}\,\frac{\partial U_{in}}{\partial \delta_{i}}\right)\left(\delta_{in}\,\frac{\partial U_{in}}{\partial t_{j,i}}\right)=o(1),\nonumber\\
\sigma_{kj}:=&\displaystyle\int_{\Omega}\nabla\left(\delta_{in}\,\frac{\partial U_{in}}{\partial t_{k,i}}\right)\nabla\left(\delta_{in}\,\frac{\partial U_{in}}{\partial t_{j,i}}\right)+\lambda_n\displaystyle\int_{\Omega}\left(\delta_{in}\,\frac{\partial U_{in}}{\partial t_{k,i}}\right)\left(\delta_{in}\,\frac{\partial U_{in}}{\partial t_{j,i}}\right)=o(1),
\end{align}
for some positive constants $\widetilde{A_0},\,\widetilde{A_1}$. On the other hand, from \eqref{new-psi}, we have 
\begin{align}
    \label{w-2}
  \langle w_{in},w_{in}\rangle=& \langle \psi_{in},w_{in}\rangle-\langle i_{\lambda_n}^{*}\left(3V_{in}^2\psi_{in}-\beta\left(V_{jn}^2\psi_{in}+2V_{in}V_{jn}\psi_{jn}\right)\right),w_{in}\rangle-\langle h_{in},w_{in}\rangle\nonumber\\
  =&\sum_{k=0}^3c_{in}^{(k)}\Bigg[\displaystyle\int_{\Omega}\left(\lambda_n Z_{k,in}\psi_{in}+6U_{in}W_{in}Z_{k,in}\psi_{in}-3W^2_{in}Z_{k,in}\psi_{in}\right)+\displaystyle\int_{\partial\Omega}\frac{\partial Z_{k,in}}{\partial \nu}\psi_{in}\nonumber\\
  &\qquad\,\quad\,\qquad-\beta\displaystyle\int_{\Omega}Z_{k,in}\left(V_{jn}^2\psi_{in}+2V_{in}V_{jn}\psi_{jn}\right)\Bigg]\nonumber\\
  =&\sum_{k=0}^3c_{in}^{(k)}\left(\lambda_n\delta_{in}|\ln{\delta_{in}}|^{\frac{1}{2}}+\delta_{in}+\delta^{\frac{3}{2}}_{in}|\ln{\delta_{in}}|^{\frac{1}{2}}+\left(\displaystyle\int_{\partial\Omega}\left|\frac{\partial Z_{k,in}}{\partial \nu}\right|^{\frac{3}{2}}\right)^{\frac{2}{3}}+{|\beta|\delta_{in}\delta_{jn}}\right)\|\bm{\psi_{n}}\|_{H}\nonumber\\ 
  \longrightarrow&\, 0,
\end{align}
since \begin{align*}
\left(\displaystyle\int_{\partial\Omega}\left|\frac{\partial Z_{k,in}}{\partial \nu}\right|^{\frac{3}{2}}\right)^{\frac{2}{3}}=&\begin{cases}
\delta_{in}\left(\displaystyle\int_{\partial\Omega}\left|\frac{\partial^{2} U_{in}}{\partial\nu\partial \delta_i}\right|^{\frac{3}{2}}\right)^{\frac{2}{3}}\lesssim \delta_{in}|\ln{\delta_{in}}|^{\frac{2}{3}},\quad &k=0,\\
\delta_{in}\left(\displaystyle\int_{\partial\Omega}\left|\frac{\partial^{2} U_{in}}{\partial\nu\partial t_{k,i}}\right|^{\frac{3}{2}}\right)^{\frac{2}{3}}\lesssim \delta_{in},\quad &k=1,2,3.
    \end{cases}
\end{align*}
Thus, from \eqref{w-1} and \eqref{w-2}, we conclude that $w_{{in}} \rightarrow 0$ strongly in $H^{1}(\Omega)$ as $n \rightarrow+\infty$, which implies that $\bm{w_{{n}}} \rightarrow 0$ strongly in $H$ as $n \rightarrow+\infty$.
\medskip

{\textbf{Step 2.}~We prove that  for $i=1,2$, \( \widetilde{\psi}_{in}\rightharpoonup 0 \) weakly in \(D^{1,2}(\R^4_{+})\) as $n \rightarrow+\infty$, where for any $x\in\Omega$,
$$
    \label{trans}\widetilde{\psi}_{in}(y):=\delta_{in}\psi_{in}\left(\delta_{in} y+\xi_{in}\right),\quad y\in\Omega_{in}:=\frac{\Omega-\xi_{in}}{\delta_{in}},\quad i=1,2.
$$ }
On the one hand, define \begin{align}
    \label{def-vn}v_{in}:=\psi_{in}-h_{in}-w_{in},
\end{align} then by the definition of $i^*_{\lambda}$, \eqref{new-psi} is equivalent to
\begin{align}\label{new-psii}
    \begin{cases}
-\Delta v_{in}+\lambda_n  v_{in}- 3V_{in}^2\,v_{in}=3V_{in}^2( h_{in}+w_{in})-\beta\left(V_{jn}^2\psi_{in}+2V_{in}V_{jn}\psi_{jn}\right), &\text{ in $\Omega$},\\
\partial_{\nu}v_{in}=0,&\text{ on $\partial\Omega$.}
\end{cases}
\end{align}
From the definition of weak solutions to \eqref{new-psii}, we derive that for any $\Phi\in H^{1}(\Omega)$, 
\begin{align}
    \label{weak-solu}
    &\displaystyle\int_{\Omega}\nabla v_{in}\nabla\Phi+\lambda_n\displaystyle\int_{\Omega}v_{in}\,\Phi -3\displaystyle\int_{\Omega}V_{in}^2\,v_{in}\,\Phi\nonumber\\=&3\displaystyle\int_{\Omega}V_{in}^2( h_{in}+w_{in})\Phi-\beta\displaystyle\int_{\Omega}\left(V_{jn}^2\psi_{in}+2V_{in}V_{jn}\psi_{jn}\right)\Phi.
\end{align}
Now  choose $\chi\in D^{1,2}(\R^4_+)$ and define $$\widehat{\chi_{in}}(x):=\frac{1}{\delta_{in}}\chi\left(\frac{x-\xi_{in}}{\delta_{in}}\right)\in H^{1}(\Omega),\quad i=1,2.$$
Replacing $\Phi$ in \eqref{weak-solu} by $\widehat{\chi_{in}}$, we have
\begin{align}
    \label{weak-solu-22}
   &\displaystyle\int_{\Omega}\nabla v_{in}\nabla\widehat{\chi_{in}}+\lambda_n\displaystyle\int_{\Omega}v_{in}\,\widehat{\chi_{in}} -3\displaystyle\int_{\Omega}V_{in}^2\,v_{in}\,\widehat{\chi_{in}}\nonumber\\=&3\displaystyle\int_{\Omega}V_{in}^2( h_{in}+w_{in})\widehat{\chi_{in}}-\beta\displaystyle\int_{\Omega}\left(V_{jn}^2\psi_{in}+2V_{in}V_{jn}\psi_{jn}\right)\widehat{\chi_{in}}.
\end{align}

In addition, for any $x\in\Omega$, define
\begin{align}
    \label{trans-2}\widetilde{v}_{in}(y):=\delta_{in}v_{in}\left(\delta_{in} y+\xi_{in}\right),\quad y\in\Omega_{in},
\end{align} 
 we claim that $\widetilde{v}_{in}\in H^1\left(\Omega_{in}\right)$. Indeed, a direct computation yields that $$\displaystyle\int_{\Omega_{in}}|\nabla\widetilde{\psi}_{in}|^2\leq \displaystyle\int_{\Omega_{in}}|\nabla\widetilde{\psi}_{in}|^2+\lambda_n\delta_{in}^2\displaystyle\int_{\Omega_{in}}|\widetilde{\psi}_{in}|^2\lesssim\|\bm{\psi_n}\|^2_{H}=1,$$
which implies that $\widetilde{\psi}_{in}\in H^1\left(\Omega_{in}\right)$. Then combining with $h_{in},\,w_{in}\to 0$ strongly in $H^1(\Omega)$, we obtain that $\widetilde{v}_{in}\in H^1\left(\Omega_{in}\right)$. Thus, there exists $\widetilde{v}_{i}\in D^{1,2}(\R^4_+)$ such that $\widetilde{v}_{in}\rightharpoonup \widetilde{v}_i$ weakly in $D^{1,2}(\R^4_+)$.
Moreover, from \eqref{weak-solu-22}, we derive that for any   $\beta\in(0,b)$,
\begin{align*}
    &\displaystyle\int_{\Omega_{in}}\nabla \widetilde{v}_{in}\nabla{\chi}+\lambda_n\delta_{in}^2\displaystyle\int_{\Omega_n}\widetilde{v}_{in}\,{\chi} \nonumber\\&-3\displaystyle\int_{\Omega_{in}}\Big(U_{0,1}^2-2\delta_{in}U_{0,1}\,W_{in}\left(\delta_{in} y+\xi_{in}\right)+\delta_{in}^2W_{in}^2\left(\delta_{in} y+\xi_{in}\right)\Big)\widetilde{v}_{in}\,{\chi}\nonumber\\=&3\displaystyle\int_{\Omega_{in}}{\chi}\left(\widetilde{h}_{in}+\widetilde{w}_{in}\right)\Big(U_{0,1}^2-2\delta_{in}U_{0,1}\,W_{in}\left(\delta_{in} y+\xi_{in}\right)+\delta_{in}^2W_{in}^2\left(\delta_{in} y+\xi_{in}\right)\Big)\nonumber\\&+2\beta\,{\delta}^3_{in}\displaystyle\int_{\Omega_{in}}V_{in}\left(\delta_{in} y+\xi_{in}\right)V_{jn}\left(\delta_{in} y+\xi_{in}\right)\psi_{jn}\left(\delta_{in} y+\xi_{in}\right){\chi}\nonumber\\&-\beta\,{\delta}^2_{in}\displaystyle\int_{\Omega_{in}}V^2_{jn}\left(\delta_{in} y+\xi_{in}\right)\widetilde{\psi}_{in}(y){\chi}\nonumber\\=&o_n(1)+|\beta|O\left(\delta_{in}\delta_{jn}|\ln{\delta_{in}}|+\delta_{in}^2\delta_{jn}^2\right)\longrightarrow 0,\qquad\text{as $n\to+\infty$,}
\end{align*}
then we have
\begin{align*}
    \displaystyle\int_{\R^4_+}\nabla \widetilde{v}_i\,\nabla{\chi}-3\displaystyle\int_{\R^4_+}U^2_{0,1}\widetilde{v}_i\,{\chi}=0,\quad\text{for any $\chi\in D^{1,2}(\R^4_+)$,}
\end{align*}
which implies  that $\widetilde{v}_i$ solves
\begin{align*}
    \begin{cases}
-\Delta \widetilde{v}_i= 3U_{0,1}^2\,\widetilde{v}_i, &\text{ in $\R^4_+$},\\
\partial_{\nu} \widetilde{v}_i=0,&\text{ on $\partial\R^4_+$,}
\end{cases}
\end{align*}
and { $$\widetilde{v}_i \mathop{//}  span\left\{\frac{\partial U_{0,1}}{\partial s_j},\,j=0,1,2,3\right\},$$}
where $\frac{\partial U_{0,1}}{\partial s_0}:=\frac{\partial U_{0,1}}{\partial \delta}$, and $\frac{\partial U_{0,1}}{\partial s_j}:=\frac{\partial U_{0,1}}{\partial t_j},\,j=1,2,3.$ In particular,
$$\frac{\partial U_{0,1}}{\partial s_j}\lesssim U_{0,1},\quad j=0,1,2,3.$$

\medskip
On the other hand, we claim that for $i=1,2$,
$$\left\langle\widetilde{v}_i,\frac{\partial U_{0,1}}{\partial s_j}\right\rangle=0,\quad j=0,1,2,3.$$
Since $\psi_{in}\in K_i^{\perp}$, that is,
\begin{align*}
    \displaystyle\int_{ \Omega} \nabla \psi_{in} \cdot \nabla Z_{j,in}+\lambda_n\displaystyle\int_{ \Omega}\psi_{in}\,Z_{j,in}=0,\quad j=0,1,2,3.
\end{align*}
then by \eqref{def-vn} and the fact that $h_{in},\,w_{in}\to 0$ strongly in $H^1(\Omega)$, we have
\begin{align*}
    &\displaystyle\int_{ \Omega} \nabla v_{in} \cdot \nabla Z_{j,in}+\lambda_n\displaystyle\int_{ \Omega}v_{in}\,Z_{j,in}\nonumber\\=&-\left(\displaystyle\int_{ \Omega} \nabla \left(h_{in}+w_{in}\right) \cdot \nabla Z_{j,in}+\lambda_n\displaystyle\int_{ \Omega}\left(h_{in}+w_{in}\right) \,Z_{j,in}\right)\nonumber\\
    \longrightarrow&0,\quad\text{as $n\to+\infty$.}
\end{align*}
Then taking the transformation \eqref{trans-2}, we have
\begin{align*}
    &\displaystyle\int_{ \Omega_{in}} \nabla \widetilde{v}_{in} \cdot \nabla \left(\frac{\partial U_{0,1}}{\partial s_j}\right)+\lambda_n\delta_{in}^2\displaystyle\int_{\Omega_{in}}\widetilde{v}_{in}\frac{\partial U_{0,1}}{\partial s_j}\nonumber\\=&-\left(\displaystyle\int_{ \Omega_{in}} \nabla \left(\widetilde{h}_{in}+\widetilde{w}_{in}\right) \cdot\nabla \left(\frac{\partial U_{0,1}}{\partial s_j}\right)+\lambda_n\delta_{in}^2\displaystyle\int_{ \Omega_{in}}\left(\widetilde{h}_{in}+\widetilde{w}_{in}\right) \, \left(\frac{\partial U_{0,1}}{\partial s_j}\right)\right)\nonumber\\
    \longrightarrow&0,\quad\text{as $n\to+\infty$.}
\end{align*}
Passing to the limit $n\to+\infty$, we derive that
\begin{align}
    \label{final} \displaystyle\int_{ \R^4_+} \nabla \widetilde{v}_i \cdot \nabla \left(\frac{\partial U_{0,1}}{\partial s_j}\right)=0.
\end{align}

Thus, from \eqref{final}, we conclude that $ \widetilde{v}_i\equiv 0$, which implies that 
\begin{align*}
    \widetilde{\psi}_i=\widetilde{v}_i+\widetilde{h}_i+\widetilde{w}_i=0,\quad\text{where $\widetilde{h}_i=\lim_{n\to+\infty}\widetilde{h}_{in}=0$ and $\widetilde{w}_i=\lim_{n\to+\infty}\widetilde{w}_{in}=0$.} 
\end{align*}
Therefore, for $i=1,2$, \( \widetilde{\psi}_{in}\rightharpoonup 0 \) weakly in \(D^{1,2}(\R^4_{+})\) as $n \rightarrow+\infty$.

{\textbf{Step 3.}~We prove that \( {\psi}_{in}\to 0 \) strongly in \(H^{1}(\Omega)\) as $n \rightarrow+\infty$}. For this aim, we firstly prove that  \( {v}_{in}\to 0 \) strongly in \(H^{1}(\Omega)\) as $n \rightarrow+\infty$.  

Indeed, a direct calculation yields that
\begin{align*}
    \|v_{in}\|^2=&\underbrace{3\displaystyle\int_{ \Omega} \left(U_{in}^2-2U_{in}W_{\lambda_n,in}+W_{\lambda_n,in}^2\right)\,v_{in}^2}_{:=(i)}+\underbrace{3\displaystyle\int_{ \Omega} V_{in}^2\,v_{in}(h_{in}+w_{in})}_{:=(ii)}\\
    &-\beta\displaystyle\int_{ \Omega} \underbrace{\big(V_{jn}^2v_{in}^2}_{:=(iii)}-\underbrace{V_{jn}^2v_{in}\left(h_{in}+w_{in}\right)}_{:=(iv)}\big)\\
    &+2\beta\displaystyle\int_{ \Omega} \big(\underbrace{V_{in}\,V_{jn}\,v_{in}\,v_{jn}}_{:=(v)}-\underbrace{V_{in}\,V_{jn}\,v_{in}\,(h_{jn}+w_{jn})}_{:=(vi)}\big),
\end{align*}
where
\begin{align*}
    (i)
    =&\underbrace{3\displaystyle\int_{ \Omega_{in}} U_{0,1}^2\,\widetilde{v}_{{in}}^2}_{:=(1)}+\underbrace{O\left(\left\|W_{\lambda_n,in}\right\|_{L^4(\Omega)}^2\,\left\|v_{in}\right\|_{L^4(\Omega)}^2+\left\|U_{in}\right\|_{L^4(\Omega)}\,\left\|W_{\lambda_n,in}\right\|_{L^4(\Omega)}\,\left\|v_{in}\right\|_{L^4(\Omega)}^2\right)}_{:=(2)},
\end{align*}
since $U_{0,1}^2\in L^{2}(\R^4_{+})$ and \( \widetilde{v}_{in}\rightharpoonup 0 \) weakly in \(D^{1,2}(\R^4_{+})\) as $n \rightarrow+\infty$, then $(1)\to 0$ as $n \rightarrow+\infty$. Moreover, by \eqref{properties-1}, we derive that as $n \rightarrow+\infty$,
\begin{align}\label{W-N}
\left\|W_{\lambda_n,in}\right\|_{L^4(\Omega)}^4=&\lambda_n^4\delta_{in}^4\left(\displaystyle\int_{ \Omega}\left|W\left(\sqrt{\lambda_n}|x-\xi_n|\right)\right|^4\right)
\lesssim\lambda^2_{n}\delta_{in}^4\left(\ln{\lambda_n}\right)^4\to 0,
\end{align}
then  $$(2)= o_{n}(1)\left\|v_{in}\right\|_{H^1(\Omega)}^2\to 0,\quad\text{as $n \rightarrow+\infty$,}$$  which implies that \begin{align}
    \label{(i)}
    (i)\to 0,\quad\text{as $n \rightarrow+\infty$.}
\end{align}
In addition, since $h_{in},\,w_{in}\to 0$  strongly in \(H^{1}(\Omega)\) as $n \rightarrow+\infty$, then
\begin{align}
    \label{(ii)}
    (ii)\lesssim\left\|V_{in}\right\|^2_{L^4(\Omega)}\,\left\|h_{in}+w_{in}\right\|_{L^4(\Omega)}\,\left\|v_{in}\right\|_{L^4(\Omega)}\to 0,\quad\text{as $n \rightarrow+\infty$.}
\end{align}
Moreover, since $\beta>0$, then 
\begin{align}
    \label{new-iii}
    (iii)\leq 0,
\end{align}
and for any $\beta\in (0,b)$, 
\begin{align}\label{new-iv}
    (iv)\lesssim|\beta|\left\|V_{jn}\right\|^2_{L^4(\Omega)}\,\left\|h_{in}+w_{in}\right\|_{L^4(\Omega)}\,\left\|v_{in}\right\|_{L^4(\Omega)}\to 0,\quad\text{as $n \rightarrow+\infty$.}
\end{align}
In addition, 
\begin{align}
    \label{(iii)}
    (v)+(vi)\lesssim&|\beta|\left\|V_{in}V_{jn}\right\|_{L^2(\Omega)}\left\|v_{in}\right\|_{L^4(\Omega)}\left(\left\|v_{jn}\right\|_{L^4(\Omega)}+\left\|h_{jn}+w_{jn}\right\|_{L^4(\Omega)}\right)\nonumber\\\lesssim&|\beta|\Big(\delta_{1n}\delta_{2n}|\ln{\left(\delta_{1n}\delta_{2n}\right)}|^{\frac{1}{2}}\Big)\left(\left\|v_{jn}\right\|_{L^4(\Omega)}+o_n(1)\right)\left\|v_{in}\right\|_{L^4(\Omega)}\to 0,\quad\text{as $n \rightarrow+\infty$.}
\end{align}



Thus, from \eqref{(i)}, \eqref{(ii)}, \eqref{new-iii}, \eqref{new-iv} and  \eqref{(iii)}, we conclude that 
\begin{align*}
    \left\|v_{in}\right\|_{H^1(\Omega)}^2\leq 0+o_{n}(1)\to 0,\quad\text{as $n \rightarrow+\infty$},
\end{align*}
which implies that \( {v}_{in}\to 0 \) strongly in \(H^{1}(\Omega)\) as $n \rightarrow+\infty$. Then since $h_{in},\,w_{in}\to 0$  strongly in \(H^{1}(\Omega)\) as $n \rightarrow+\infty$, we obtain that
\( {\psi}_{in}\to 0 \) strongly in \(H^{1}(\Omega)\) as $n \rightarrow+\infty$, $i=1,2$. 

\medskip
\textbf{Step 4.}~We prove that the operator $$\Pi^{\perp}_{\bm{d},\bm{\xi}}\circ\bm{\mathcal{L}}_{\bm{d},\bm{\xi}}:= \left(\Pi_1^{\perp}\circ\mathcal{L}_1,\Pi_2^{\perp}\circ\mathcal{L}_2\right)$$ is an invertible operator. Since $i_\lambda^*$ is a compact operator from $L^{\frac{4}{3}}(\Omega)$ to $H^1(\Omega)$, and $\Pi_i^{\perp}\circ\mathcal{L}_i:=Id-K_i$, where $K_i$ is a compact operator. Since we have proved that 
$$\left\|\Pi^{\perp}_{\bm{d},\bm{\xi}}\circ\bm{\mathcal{L}}_{\bm{d},\bm{\xi}}(\bm{\psi})\right\|_{H}=\sum_{i=1}^2\left\|\Pi_i^{\perp}\circ\mathcal{L}_i(\bm{\psi})\right\|\geq C_0\left\|\bm{\psi}\right\|_{H}, \quad\text{$\forall \bm{\psi}\in \bm{K^{\perp}}_{\bm{d},\bm{\xi}}$,}$$
 for some constant $C_0>0$, which implies that $\Pi^{\perp}_{\bm{d},\bm{\xi}}\circ\bm{\mathcal{L}}_{\bm{d},\bm{\xi}}$ is injective. Thus, by Fredholm's alternative Theorem, $\Pi^{\perp}_{\bm{d},\bm{\xi}}\circ\bm{\mathcal{L}}_{\bm{d},\bm{\xi}}$ is surjective. Hence, the operator $\Pi^{\perp}_{\bm{d},\bm{\xi}}\circ\bm{\mathcal{L}}_{\bm{d},\bm{\xi}}$ is invertible, and $\left(\Pi^{\perp}_{\bm{d},\bm{\xi}}\circ\bm{\mathcal{L}}_{\bm{d},\bm{\xi}}\right)^{-1}$ is continuous. This completes the proof.
 
    \end{proof}

\end{document}
